\setcounter{section}{2}

\Needspace{5\baselineskip}
\hypertarget{state-distributions-and-occupancy-measures}{%
\section{State Distributions and Occupancy Measures}\label{state-distributions-and-occupancy-measures}}

\Needspace{5\baselineskip}
\hypertarget{i.-definitions-of-state-distributions-and-occupancy-measures}{%
\subsection{I. Definitions of State Distributions and Occupancy Measures}\label{i.-definitions-of-state-distributions-and-occupancy-measures}}

\Needspace{5\baselineskip}
The state distribution is defined as the discounted probability that the agent occupies state \(s\):

\Needspace{5\baselineskip}
Define \(\nu^\pi(s)\) as the discounted state-visitation distribution under policy \(\pi\):

\[
\nu^\pi(s) = (1-\gamma)\sum_{t=0}^{\infty}\gamma^t P_t^\pi(s)
\]

\(P_t^\pi(s)\): the probability that the agent occupies state \(s\) at time \(t\).

\(\gamma\): the discount factor, \(0 \le \gamma < 1\).

\(1-\gamma\): the normalization factor.

\Needspace{5\baselineskip}
We can then define the occupancy measure, the discounted probability of the agent taking action a in state s:

\[
\rho^\pi(s,a) = (1-\gamma)\sum_{t=0}^{\infty}\gamma^t P_t^\pi(s)\pi(a|s)
\]

\Needspace{5\baselineskip}
This gives the relationship:

\[
\rho^\pi(s,a) = \nu^\pi(s)\pi(a|s)
\]

\Needspace{5\baselineskip}
\hypertarget{ii.-derivation}{%
\subsection{II. Derivation}\label{ii.-derivation}}

\begin{enumerate}
\def\labelenumi{\arabic{enumi}.}
\tightlist
\item
  Why is a discount factor needed?
\end{enumerate}

\Needspace{5\baselineskip}
We usually seek to maximize cumulative discounted return:

\[
J(\pi) = \mathbb{E}_\pi\left[\sum_{t=0}^{\infty}\gamma^t R(s_t,a_t)\right]
\]

\Needspace{5\baselineskip}
By linearity of expectation, move the sum outside the expectation:

\[
J(\pi) = \sum_{t=0}^{\infty}\gamma^t\mathbb{E}_\pi[R(s_t,a_t)]
\]

\[
J(\pi) = \sum_{t=0}^{\infty}\gamma^t\sum_{s,a}P_t^\pi(s)\pi(a|s)R(s,a)
\]

\Needspace{5\baselineskip}
Now move the sum over \(s\) and \(a\) to the front:

\[
J(\pi) = \sum_{s,a}R(s,a)\left(\sum_{t=0}^{\infty}\gamma^t P_t^\pi(s)\pi(a|s)\right)
\]

\Needspace{5\baselineskip}
The parenthesized term is the unnormalized occupancy measure. To turn it into a probability distribution (summing to 1), multiply and divide by \((1-\gamma)\):

\[
J(\pi) = \frac{1}{1-\gamma}\sum_{s,a}R(s,a)\left((1-\gamma)\sum_{t=0}^{\infty}\gamma^t P_t^\pi(s)\pi(a|s)\right)
\]

\Needspace{5\baselineskip}
The parenthesized part is \(\rho^\pi(s,a)\), so:

\[
J(\pi) = \frac{1}{1-\gamma}\sum_{s,a}\rho^\pi(s,a)R(s,a) = \frac{1}{1-\gamma}\mathbb{E}_{(s,a)\sim\rho^\pi}[R(s,a)]
\]

In reinforcement learning, the goal is often to find a good policy. A policy can be represented by a probability distribution, and the optimization objective is precisely to maximize the expectation of return \(R\) under \(\rho^\pi\). The problem thereby becomes finding such a distribution. A state-action pair's probability is its contribution weight to total value, consistent with our understanding of Markov decision processes.

\begin{enumerate}
\def\labelenumi{\arabic{enumi}.}
\setcounter{enumi}{1}
\tightlist
\item
  Why is a normalization factor needed?
\end{enumerate}

This is a crucial mathematical detail. We must prove that \(\nu^\pi(s)\) is a valid probability distribution, summing to 1 over all states.

\Needspace{5\baselineskip}
Sum over every state \(s\):

\[
\sum_s \nu^\pi(s) = \sum_s\left[(1-\gamma)\sum_{t=0}^{\infty}\gamma^t P_t^\pi(s)\right]
\]

\Needspace{5\baselineskip}
Interchange the sums (assuming convergence):

\[
= (1-\gamma)\sum_{t=0}^{\infty}\gamma^t\left(\sum_s P_t^\pi(s)\right)
\]

\Needspace{5\baselineskip}
Whatever \(t\) is, the agent must occupy some state, so \(\sum_s P_t^\pi(s)=1\). The expression becomes:

\[
= (1-\gamma)\sum_{t=0}^{\infty}\gamma^t \cdot 1 = 1
\]

\Needspace{5\baselineskip}
\hypertarget{iii.-mathematical-meaning-of-occupancy-measures-sampling-time-geometrically}{%
\subsection{III. Mathematical Meaning of Occupancy Measures: Sampling Time Geometrically}\label{iii.-mathematical-meaning-of-occupancy-measures-sampling-time-geometrically}}

We usually understand ``visitation probability'' as the percentage of all time steps spent in a state over the long run (a simple frequency statistic).

With discount factor \(\gamma\), however, \(\rho^\pi(s,a)\) is clearly no longer a simple ``time frequency,'' because it emphasizes earlier times over later ones.

It is still called a ``probability'' because we have changed the probability distribution used to ``sample time \(t\).''

Traditional view (undiscounted, \(\gamma=1\))

Every time \(t\) is assumed equally likely to be sampled (a mathematically uniform distribution is actually impossible over infinite time, so a limit is usually taken). This leads to the familiar ``stationary distribution,'' simply the time-average proportion.

Discounted view (discounted, \(\gamma<1\))

\(\rho^\pi(s,a)\) actually defines a new random experiment: instead of sampling time uniformly, sample \(t\) according to a geometric distribution.

\Needspace{5\baselineskip}
Sampling rules:

\begin{itemize}
\tightlist
\item
  The probability of \(t=0\) is \((1-\gamma)\).
\item
  The probability of \(t=1\) is \((1-\gamma)\gamma\).
\item
  The probability of \(t=2\) is \((1-\gamma)\gamma^2\).
\item
  \ldots{}
\item
  The probability of time \(t\) is \(P(T=t)=(1-\gamma)\gamma^t\), where \(T\) is the sampled time.
\end{itemize}

\Needspace{5\baselineskip}
Its meaning is now:

If the observation time is uncertain and an instant is drawn randomly according to this ``sooner matters more'' rule, what is the probability that the agent is at \((s,a)\) at that instant?

\Needspace{5\baselineskip}
Another interpretation is:

\Needspace{5\baselineskip}
For greater intuition, introduce random termination (also the classic physical interpretation of \(\gamma\)):

\Needspace{5\baselineskip}
Suppose the game does not continue indefinitely. Instead, at the end of each step, God rolls a die:

\begin{itemize}
\tightlist
\item
  With probability \(\gamma\), the game continues.
\item
  With probability \(1-\gamma\), the game suddenly ends.
\end{itemize}

\Needspace{5\baselineskip}
Then \(\rho^\pi(s,a)\) measures:

The probability that the agent is in state \(s\) and taking action \(a\) in the game-over snapshot, at the instant the game ``suddenly ends''.

\Needspace{5\baselineskip}
This interpretation matches the formula perfectly:

\begin{itemize}
\tightlist
\item
  The probability that the game ends exactly at \(t\) is \((1-\gamma)\gamma^t\).
\item
  The probability that the agent is at \((s,a)\) at \(t\) is \(P_t^\pi(s,a)\).
\item
  Summing over all possible termination times gives the total probability of observing \((s,a)\) at termination.
\end{itemize}

Because the game must end at some point, these probabilities sum to 1 over all \((s,a)\). It is therefore a proper probability distribution.

\Needspace{5\baselineskip}
\hypertarget{iv.-recursive-representation-of-the-state-distribution}{%
\subsection{IV. Recursive Representation of the State Distribution}\label{iv.-recursive-representation-of-the-state-distribution}}

When calculating \(\nu^\pi(s')\), we sum the weighted occurrences of the agent at \(s'\) across all time steps from \(t=0\) to \(t=\infty\).

\Needspace{5\baselineskip}
Every occurrence has one of two possible timings:

\begin{enumerate}
\def\labelenumi{\arabic{enumi}.}
\tightlist
\item
  It occurs at \(t=0\): the ``initial state.''
\item
  It occurs at \(t>0\): a ``transition from the previous time step.''
\end{enumerate}

Total probability at \(s'\) = weight of starting there + weight of transitions from elsewhere.

\[
\nu^\pi(s') = (1-\gamma)\nu_0(s') + \gamma\int P(s'|s,a)\pi(a|s)\nu^\pi(s)\,ds\,da
\]

The first term \((1-\gamma)\nu_0(s')\) is the initial contribution, and the second \(\gamma\int P(s'|s,a)\pi(a|s)\nu^\pi(s)\,ds\,da\) is the transition contribution.

Here \(\nu_0(s')\) is the probability \(P(s_0=s')\) of initially being at \(s'\), and is not the same function as \(\nu^\pi\).

First term: \((1-\gamma)\nu_0(s')\)

\begin{itemize}
\tightlist
\item
  Meaning: the weighted probability of being at \(s'\) at time \(t=0\).
\item
  \(\nu_0(s')\): the initial-state distribution, \(P(s_0=s')\).
\item
  \((1-\gamma)\): a normalization coefficient for time weights.
\item
  Recall the definition of \(\nu^\pi\): \(\nu^\pi(s)=(1-\gamma)\sum_{t=0}^{\infty}\gamma^t P(s_t=s)\).
\item
  When \(t=0\), \(\gamma^0=1\).
\item
  The contribution of time \(t=0\) is therefore coefficient \((1-\gamma)\) multiplied by probability \(\nu_0(s')\).
\end{itemize}

Second term: \(\gamma\int P(s'|s,a)\pi(a|s)\nu^\pi(s)\,ds\,da\)

\begin{itemize}
\tightlist
\item
  Meaning: the weighted probability of reaching \(s'\) in one transition from any state \(s\) at the previous time.
\item
  \(\nu^\pi(s)\): the probability density of state \(s\) at the previous time.
\item
  \(\pi(a|s)\): the probability of choosing action \(a\) in \(s\).
\item
  \(P(s'|s,a)\): the probability of moving from \(s\) to \(s'\) after taking \(a\).
\item
  \(\int \cdots dsda\): integrates over all possible ``predecessor'' states and actions (or sums \(\sum\) in discrete spaces), collecting all possibilities leading to \(s'\).
\item
  \(\gamma\): the crucial discount factor.
\item
  Moving from \(s\) to \(s'\) takes one time step.
\item
  If the previous time (say \(t\)) has weight \(\gamma^t\), arrival at \(s'\) occurs at \(t+1\), with weight \(\gamma^{t+1}\).
\item
  The weight is discounted by a factor of \(\gamma\) relative to the previous time, so multiplication by \(\gamma\) is required.
\end{itemize}

\Needspace{5\baselineskip}
\hypertarget{v.-theorems-concerning-occupancy-measures}{%
\subsection{V. Theorems Concerning Occupancy Measures}\label{v.-theorems-concerning-occupancy-measures}}

Two further theorems follow.

\Needspace{5\baselineskip}
Theorem 1: when an agent interacts with the same MDP under policies \(\pi_1\) and \(\pi_2\), the resulting occupancy measures \(\rho^{\pi_1}\) and \(\rho^{\pi_2}\) satisfy:

\[
\rho^{\pi_1} = \rho^{\pi_2} \Longleftrightarrow \pi_1 = \pi_2
\]

\Needspace{5\baselineskip}
Theorem 2: given a valid occupancy measure \(\rho\), the unique policy that generates it is:

\[
\pi_\rho = \frac{\rho(s,a)}{\sum_{a'}\rho(s,a')}
\]

\FloatBarrier
\Needspace{5\baselineskip}
\hypertarget{references-44}{%
\subsection{References}\label{references-44}}

\vspace{0.25em}
\begin{itemize}
\tightlist
\item
  \href{https://hrl.boyuai.com/}{Hands-on Reinforcement Learning (translated title; in Chinese)}.
\end{itemize}

\clearpage
